
\def\PUDcodigo{EME131}

\if\COMPILAR0
    \def\PUDcodacad{04507.60}
\else
    \def\PUDcodacad{04506.60}
\fi

\def\PUDdisciplina{Equações Diferenciais Ordinárias}

\def\PUDsemestre{Opt}

\def\PUDhoras{80}

%NOVOS ---------------------------------------------------
\def\PUDhorasTeoria{80}
\def\PUDhorasPratica{0}

\def\PUDinterECA{---}

\def\PUDatuaisECA{---}

\def\PUDrecursos{Projetor multimídia; Quadro branco e pincel.}

%----------------------------------------------------------

\def\PUDcreditos{4}

\def\PUDprerequisito{ \hyperref[PUD:ACE006]{ACE006} }

\if\COMPILAR0
   \def\PUDpreacad{ \hyperref[PUD:ACE006]{Cálculo 2 (04507.7)} }
\else
   \def\PUDpreacad{ \hyperref[PUD:ACE006]{Cálculo 2 (04506.7)} }
\fi

\def\PUDobjetivos{Familiarizar-se com os conceitos e definições de Equações Diferenciais Ordinárias, de forma a aplicá-los na respectiva área de atuação e nas disciplinas que envolvam a matemática aplicada. Reconhecer a importância e a influência que a matemática exerce no cotidiano e no progresso de pesquisas científicas. Especificamente, seria: Desenvolver o conceito de equações diferenciais ordinárias (EDO’s); Estudar técnicas de resolução de EDO’s; Estudar alguns exemplos de aplicação, reconhecendo a importância da disciplina em outras áreas. }

\def\PUDementa{Equações diferenciais de 1ª ordem; Campo Vetorial; Equações Diferenciais Separáveis; Equações Diferencias Lineares de 1ª ordem e o fator integrante. Equações Diferencias Lineares de 2ª ordem; EDL homogêneas; Princípio da superposição; Transformada de Laplace; Resolução de EDO's utilizando transformada de Laplace.}

\def\PUDconteudo{ & Equações Diferenciais Ordinárias de 1ª Ordem: Definição e exemplos de equação diferencial; Equações Diferenciais Ordinárias lineares e não lineares; Equações de Variáveis Separáveis, Fator Integrante; Equações Exatas e as Redutíveis a ela por meio de fator integrante; Teorema de Existência e Unicidade das Soluções; Interpretação Gráfica das Soluções. 
\\ 
 & Equações Diferenciais Ordinária de Ordem Superior: Problema de valor inicial; dependência linear e não linear; Equações Homogêneas com Coeficiente Constante; Equações Não Homogêneas; Método dos Coeficientes Indeterminados; O Método de Variação dos Parâmetros; Solução em séries de Potências de EDO'S de 2ª Ordem; Aplicações. 
\\ 
 & Sistemas de Equações Diferenciais Lineares: Sistemas Lineares; Sistemas Lineares Homogêneos com os coeficientes constantes; Sistemas não lineares; Estabilidade de sistemas; Método de Euler e Runge-Kutta para resolução de EDO's; Aplicações. 
\\ 
 & Transformada de laplace: Obtenção da transformada das funções usuais; Tabela das transformações; Resolução das equações com coeficientes constantes, através do uso das transformadas de Laplace.
 }

\def\PUDmetodo{Aulas expositórias que podem ser teóricas e/ou práticas, onde as práticas no laboratório serão marcadas no decorrer da disciplina.}

\def\PUDavalia{A avaliação será feita com aplicação de provas teóricas e/ou práticas, além da possibilidade de inclusão de trabalhos e seminários no decorrer da disciplina.}

\def\PUDbibbasica{ &  \href{www.biblioteca.ifce.edu.br/index.asp?codigo_sophia=41966}{Boyce}, E.W.; Diprima, R.C. Equações Diferenciais Elementares e Problemas de Valores de Contorno.  9ª ed.  Rio de Janeiro, RJ: LTC, 2013.  607p.  ISBN: 9788521617563. 
\\ 
 &  \href{www.biblioteca.ifce.edu.br/index.asp?codigo_sophia=63261}{Howard} Anton, Chris Rorres. Álgebra linear com aplicações.  10ª ed.  Porto Alegre, RS: Bookman, 2012.  768p.  ISBN: 9788540701694. 
\\
&  \href{www.biblioteca.ifce.edu.br/index.asp?codigo_sophia=24186}{Leithold}, Louis.  O cálculo com geometria analítica.  3ª ed.  São Paulo: Harbra, 1994.  v.2.  ISBN: 8529402065. }

\def\PUDbibcomplementar{ & \href{www.biblioteca.ifce.edu.br/index.asp?codigo_sophia=40051}{Burden}, Richard L.  Análise numérica.  São Paulo, SP: Cengage Learning, 2013.  721p.  ISBN: 9788522106011.
\\
 & \href{www.biblioteca.ifce.edu.br/index.asp?codigo_sophia=62475}{Gilat}, Amos.  MATLAB com aplicações em engenharia.  4ª ed.  Porto Alegre, RS: Bookman, 2012.  417p.  ISBN: 9788540701861.
\\
 & \href{www.biblioteca.ifce.edu.br/index.asp?codigo_sophia=24241}{Chapman}, Stephen J.  Programação em MATLAB para engenheiros.  2ª ed.  São Paulo, SP: Cengage Learning, 2010.  410p.  ISBN: 9788522107896.
\\
 & \href{www.biblioteca.ifce.edu.br/index.asp?codigo_sophia=56952}{Nagle}, R. Kent; Saff, Edwar B.  Equações Diferenciais. E-book.  8ª ed.  Pearson.  584p.  ISBN: 9788581430836.
\\
 & \href{www.biblioteca.ifce.edu.br/index.asp?codigo_sophia=39837}{Lathi}, B. P.  Sinais e sistemas lineares.  2ª ed.  Porto Alegre, RS: Bookman, 2007.  856p.  ISBN: 9788560031139. }

\def\PUDautor{David Carneiro de Souza}

\def\PUDdata{2017-03-20}

\def\PUDrevisao{2}

\def\PUDrevisor{Samuel Vieira Dias}

\def\PUDdatarevisao{2017-07-17}

\PUDcorpo